\documentclass[11pt,a4paper]{article}
\usepackage[margin=25mm]{geometry}
\usepackage[T1]{fontenc}
\usepackage{amsmath,amssymb,amsthm,booktabs}
\usepackage[hidelinks]{hyperref}
\hypersetup{pdftitle={Sudoku Island: Correctness and Completeness},pdfauthor={Sudoku Island project}}
\newtheorem{lemma}{Lemma}
\newtheorem{theorem}{Theorem}
\setlength{\parindent}{0pt}
\setlength{\parskip}{6pt}
\title{Sudoku Island Search\\\large Correctness, Completeness, and an Even-Digit Witness}
\author{Sudoku Island project}
\date{25 November 2026}

\begin{document}
\maketitle

\section{The precise problem}
Let $b\geq1$ and $n=b^2$. A board has $n$ rows, $n$ columns, and $n$
boxes, each of size $b\times b$. A pattern is a subset $S$ of its cells.
Two selected cells are adjacent if they share a side. There is no diagonal
adjacency or wrapping at the boundary.

For a given $k\in\{0,\ldots,n\}$, we require exactly $k$ cells of $S$ in
each row, each column, and each box, and require that $S$ be nonempty and
connected. Thus $|S|=nk$, and $k=0$ is excluded by definition. The standard
problem uses $b=3$, $n=9$. Its objective is to minimize positive $k$.

This is initially a binary problem: no Sudoku digits are assumed. A separate
second stage tests whether the chosen $k=4$ pattern can be the even cells of
a genuine Sudoku. No universal extension assumption is used.

\section{Search representation}
Each row is represented by a mask $m\in\{0,\ldots,2^n-1\}$ with exactly
$k$ set bits. Bit zero denotes the leftmost column. All
$\binom{n}{k}$ such masks are generated, in increasing lexicographic order
of their selected column tuples. Rows are processed from top to bottom.

After $d$ rows, the state consists of:
\begin{itemize}
\item the depth $d$ and the $n$ column counts;
\item the $b$ box counts in the current band of $b$ rows (zero at band boundaries);
\item the selected positions in the last processed row, with labels encoding
which belong to the same component in the entire processed prefix.
\end{itemize}
An empty position has label zero. Positive labels are canonicalized by first
appearance from left to right. All components of an accepted prefix meet its
last row. At depth zero the frontier and counts are all zero.

The actual row masks are retained as a witness path. For the ordinary
existence search, they are not part of the failed-state cache key.
Section~\ref{sec:complete} justifies this omission. A later uniqueness search
adds the remaining matches against excluded patterns to this key.

\clearpage
\section{Transitions and safe rejection rules}
\begin{lemma}[Quota bounds]
Suppose a proposed row leaves column count $a_c$, with $r$ rows remaining.
A necessary condition for a completion is
\[
 a_c\leq k\leq a_c+r.
\]
If its current-band box count is $q_j$ and $t$ rows remain in that band,
a necessary condition is
\[
 q_j\leq k\leq q_j+tb.
\]
At the end of a column or band, these inequalities enforce equality to $k$.
\end{lemma}
\begin{proof}
A future row contributes either zero or one cell to a particular column,
and at most $b$ cells to a particular box. Counts never decrease. The stated
bounds follow; at $r=0$ or $t=0$, the lower and upper bounds coincide.
These are necessary conditions only, which is sufficient for safe rejection.
\end{proof}

\begin{lemma}[Exact frontier update]
The function \texttt{\_advance} computes the component partition on the new
frontier exactly, provided that no old component is lost.
\end{lemma}
\begin{proof}
Create one node per old component and one per selected new cell. Join a new
cell to its old component whenever the cell directly above it was selected,
and join horizontally adjacent selected new cells. These are precisely all
new edges. An old component already summarizes every path in the old prefix.
Consequently, connected components of this small union graph are exactly
components of the enlarged prefix, restricted to its frontier. Relabeling
these components canonically changes no equivalences.
\end{proof}

\begin{lemma}[No disappearing component]
For $k>0$, a prefix transition that leaves an old component without any cell
in the new frontier cannot belong to a valid complete pattern.
\end{lemma}
\begin{proof}
Every path from an old row to a new or later row must pass through the new
frontier, since grid edges change row number by at most one. A component
that has no representative there can never acquire further connections.
The newly appended row contains $k>0$ selected cells. They lie outside that
component, so the completed pattern would be disconnected, even if this was
the final row. Rejecting the transition is therefore safe.
\end{proof}

In the implementation, the old frontier labels become union nodes. After
adding the new edges, every old label must have the same union root as some
selected new cell. Otherwise the transition is rejected. This implements the
lemma without any geometric symmetry assumption.

At depth $n$, acceptance requires exactly one positive frontier label.
All row quotas hold by construction; the quota bounds have already enforced
all column and box quotas. The frontier invariant then ensures global
connectivity, including components formed in earlier rows.

\clearpage
\section{Exhaustiveness and memoization}\label{sec:complete}
\begin{lemma}[Equivalent future problems]
Two accepted prefixes with the same depth, column counts, current-band box
counts, and canonical frontier partition admit exactly the same suffixes.
\end{lemma}
\begin{proof}
Column counts completely determine outstanding column quotas. Current-band
counts determine outstanding box quotas; earlier boxes are finished and
later boxes are untouched. Depth determines which box band is being filled.
Each future row must have exactly $k$ cells independently of the prefix.

All contacts between a suffix and a prefix pass through the last prefix row.
Its occupied positions and their connectivity partition determine which
suffix connections merge old components. Since every prefix component meets
that frontier, no other component information is needed. Thus quota tests,
frontier transitions, and final acceptance agree for every possible suffix.
\end{proof}

\begin{theorem}[Soundness, termination, and completeness]
With enough resources and no interruption, \texttt{solve} terminates and
returns SAT if and only if a valid pattern exists; otherwise it returns UNSAT.
Every returned witness satisfies the problem definition.
\end{theorem}
\begin{proof}
For $k=0$ the answer is immediately UNSAT by nonemptiness. Suppose $k>0$.
There are at most $\binom{n}{k}^{n}$ row sequences and recursion depth at
most $n+1$, so the underlying search tree is finite.

For soundness, every accepted leaf has the required quotas and one global
component by the transition invariant and final test. The implementation
also validates its witness using a separate whole-board count and graph
traversal, raising an error if that check fails.

For completeness, fix any valid full pattern. Its row masks all appear in
the candidate lists. Its prefixes satisfy the quota bounds and cannot lose
a component that must later connect to the rest. Therefore its path cannot
be rejected by these rules. A cache hit cannot remove a successful suffix:
only states for which all candidate continuations have failed are cached,
and equal keys have identical suffix problems by the preceding lemma.
Thus either the search finds another witness first, or reaches this one.
UNSAT is returned only when no valid path remains.
\end{proof}

The transition cache merely stores a pure function of the frontier and new
mask; eviction affects runtime only. The failed-state cache is local to each
quota search. A time or state limit raises an interruption that propagates
without caching the unfinished state, and produces UNKNOWN rather than UNSAT.

\begin{theorem}[Minimum search]
If quotas are tested in ascending positive order, the first SAT quota is
the minimum provided all preceding answers are UNSAT.
\end{theorem}
\begin{proof}
The SAT witness supplies an upper bound, while completeness of each earlier
UNSAT decision excludes every smaller positive quota. If a quota returns
UNKNOWN, the implementation stops without asserting a minimum.
\end{proof}

\clearpage
\section{Computed result and Sudoku completion}
The reference command is
\begin{center}
\texttt{python sudoku\_island.py --minimum --check-uniqueness}\\
\texttt{--sudoku --json docs/results.json}
\end{center}
Its abstract search gives the following reproducible state counts. The SAT
run stops at its first witness; it does not enumerate all $k=4$ solutions.
\begin{center}
\begin{tabular}{rrrr}
\toprule
$k$ & Outcome & Expanded states & Failed-state cache hits\\
\midrule
1 & UNSAT & 10 & 0\\
2 & UNSAT & 84 & 0\\
3 & UNSAT & 1,016 & 3\\
4 & SAT & 19,894 & 8,761\\
\bottomrule
\end{tabular}
\end{center}
The witness masks, from top to bottom, are
\[
(15,393,408,464,113,39,60,102,450).
\]
The separate digit search produces this complete Sudoku:
\begin{center}
\renewcommand{\arraystretch}{1.15}
\begin{tabular}{|ccc|ccc|ccc|}\hline
2&4&6&8&1&3&5&7&9\\
8&3&9&2&5&7&1&4&6\\
5&7&1&6&4&9&3&2&8\\\hline
3&9&5&7&2&1&8&6&4\\
4&1&7&9&6&8&2&3&5\\
6&2&8&5&3&4&7&9&1\\\hline
7&5&2&4&8&6&9&1&3\\
1&8&4&3&9&2&6&5&7\\
9&6&3&1&7&5&4&8&2\\\hline
\end{tabular}
\end{center}
The digit search chooses an unfilled cell of smallest current domain and
tries every admissible digit: $\{2,4,6,8\}$ on the island and
$\{1,3,5,7,9\}$ elsewhere. Used digits in the row, column and box are
excluded. Each recursive step fills a cell; each failure undoes exactly its
assignment. This finite branching search is complete for the chosen pattern.
A separate validator checks all 27 Sudoku units and exact parity agreement.

Every Sudoku with a connected even-digit set projects to a valid abstract
$k=4$ pattern. The displayed completion proves that the abstract lower bound
is attained by a genuine Sudoku. The next section establishes uniqueness of
the minimum abstract shape. Rotating and reflecting this Sudoku therefore
completes every minimum pattern. Digit assignments are not claimed to be unique.

\clearpage
\section{Uniqueness under rotations and reflections}
Let $D_4$ be the eight symmetries of the square, consisting of four rotations
and their compositions with one mirror reflection. These preserve side
adjacency, row and column quotas, and the standard box partition. No other
Sudoku permutations are used. For a reference solution $S$, let
$E=\{g(S):g\in D_4\}$, with duplicates removed.

The uniqueness check searches for a valid pattern outside $E$. At depth $d$
it carries the set $M_d$ of indices of excluded patterns whose first $d$ rows
equal the current path. Appending mask $m$ replaces this set by
\[
 M_{d+1}=\{i\in M_d : \text{row }d\text{ of }E_i=m\}.
\]
The failed-state key includes $M_d$. At a complete connected leaf, acceptance
requires $M_n=\varnothing$. Thus a pattern is rejected at the leaf exactly when
it belongs to $E$. Omitting this prefix information from the cache key would
be unsound: geometrically equivalent prefix states can have different
remaining exclusions.

\begin{theorem}[Uniqueness decision]
If the outside-orbit search returns UNSAT, every valid pattern is in $E$.
If it returns SAT, its witness and $S$ belong to different symmetry orbits.
\end{theorem}
\begin{proof}
The transition and quota arguments are unchanged. Equal augmented keys
also have identical exclusion tests for each suffix, since their matching
index sets agree. The completeness proof therefore applies to the space
of valid patterns minus $E$. UNSAT means this difference is empty. SAT gives
a validated member of the difference. UNKNOWN establishes neither conclusion.
\end{proof}

The reference $k=4$ pattern has eight distinct symmetry images, all validated
as islands. The completed search outside these images returns UNSAT after
493,104 expanded states and 434,528 failed-state cache hits (about 87 seconds
on the development machine). Hence there are exactly eight oriented
$k=4$ patterns and \textbf{one shape up to rotation and reflection}.
This counts binary islands, not Sudoku digit assignments.

Every symmetry image of the reference Sudoku remains a Sudoku, with its
even cells equal to the corresponding image of the island. In conjunction
with uniqueness this proves extendability of all minimum islands, without
a general extension assumption for other binary patterns.

\section{Evidence and scope of the guarantee}
Tests compare the witness checker with all $2^{16}=65,536$ binary $4\times4$
boards, and compare frontier transitions against direct graph traversal for
all four-column prefixes of up to three rows. Additional tests cover the
standard minimum, Sudoku completion, all twelve small-board $k=3$ solutions
via successive exclusions, symmetry orbits, and interruption handling.

The numerical minimum is computer-assisted: the proof above establishes the
algorithm, while the recorded run supplies its outcomes. Tests are supporting
evidence, not formal verification of Python or its runtime. The JSON file is
an audit record, not a separately checkable UNSAT certificate. Memory exhaustion,
external termination, and interrupted searches cannot establish UNSAT.
No claim is made that every board is explicitly visited: safe pruning and
equivalent failed states account for omitted branches.

\end{document}
